% Source-derived chapter 02: Review of hyper-K\"ahler manifolds
% Original source: riemann-roch-polynomials-hyperkahler-fourfolds
\chapter{Review of hyper-K\"ahler manifolds}
\label{riemann-roch-polynomials-hyperkahler-fourfolds-chapter-02}
\source{riemann-roch-polynomials-hyperkahler-fourfolds}

\begin{remark}[Source section]
\label{riemann-roch-polynomials-hyperkahler-fourfolds-chapter-02-source}
\source{riemann-roch-polynomials-hyperkahler-fourfolds}
\source{riemann-roch-polynomials-hyperkahler-fourfolds}
This chapter reproduces the corresponding section of the retrieved
LaTeX source, including its definitions, statements, examples, and
proofs. It has not yet been translated into Lean.
\notready
\end{remark}

% The source section title was: Review of hyper-K\"ahler manifolds
\label{sec:RRpolynomial}
\source{riemann-roch-polynomials-hyperkahler-fourfolds}



We recall in Section~\ref{subsec:BBF} the definitions of  Beauville--Bogomolov--Fujiki forms and  Fujiki constants for \hk\ manifolds.\
We then define  their Huybrechts--Riemann--Roch polynomials  and   review some of their elementary properties   in Section~\ref{subsec:RRpolynomial}.

\subsection{The Beauville--Bogomolov--Fujiki form and the Fujiki constant}\label{subsec:BBF}
\source{riemann-roch-polynomials-hyperkahler-fourfolds}

Let $X$ be a \hk\ manifold of dimension $2n$.\
There exists a canonical integral nondivisible quadratic form $q_X$ (the \emph{Beauville--Bogomolov--Fujiki form}) on $H^2(X,\Z)$ and a positive rational constant $c_X$ (the \emph{Fujiki constant}) such that
\begin{equation}\label{fuj}
\source{riemann-roch-polynomials-hyperkahler-fourfolds}
\forall \alpha\in H^2(X,\R)\qquad \int_X \alpha^{2n}=c_X\, q_X(\alpha)^n
\end{equation}
(after extending $q_X$ to a quadratic form on $H^2(X,\R)$).\
Moreover, $q_X(h)>0$ for all K\"ahler  classes~$h$ (see \cite[Proposition~23.14]{huy}).

Assume $q_X(\alpha)=0$.\
Then $\alpha^{n+1}=0$ (see for example \cite[Proposition~24.1]{huy}) and, by comparing the coefficients of $t^n$ in the relation
\[
\int_X (t\alpha+ \beta)^{2n}=c_Xq_X(t\alpha+ \beta)^n=c_X(2tq_X( \alpha, \beta)+q_X( \beta))^n,
\]
we get
\begin{equation}\label{pol}
\source{riemann-roch-polynomials-hyperkahler-fourfolds}
\forall \beta\in H^2(X,\R)\qquad \frac{1}{2^n}\binom{2n}{n}\int_X \alpha^{n}\beta^{n}= c_Xq_X( \alpha, \beta)^n.
\end{equation}
This is a particular case of the polarization of the Fujiki relation~\eqref{fuj}, which in dimension 4 takes the form
 \begin{equation}\label{eq:Fujiki3}
\source{riemann-roch-polynomials-hyperkahler-fourfolds}
 3\!\int_X \alpha_1\alpha_2\alpha_3\alpha_4 =c_X\bigl( q_X(\alpha_1,\alpha_2)q_X(\alpha_3,\alpha_4) + q_X(\alpha_1,\alpha_3)q_X(\alpha_2,\alpha_4) + q_X(\alpha_1,\alpha_4)q_X(\alpha_2,\alpha_3)\bigr)
\end{equation}
for all $\alpha_1, \alpha_2, \alpha_3,\alpha_4\in H^2(X,\Z)$.\

\subsection{The  Huybrechts--Riemann--Roch polynomial}\label{subsec:RRpolynomial}
\source{riemann-roch-polynomials-hyperkahler-fourfolds}

By \cite[Corollary~23.18]{huy}, there is a polynomial
\begin{equation*}
P_{RR,X}(T)=\sum_{i=0}^n  a_i T^i\in\Q[T]
\end{equation*}
of degree $n$  such that, for each line bundle $L$ on $X$, one has
\begin{equation}\label{rr}
\source{riemann-roch-polynomials-hyperkahler-fourfolds}
\chi(X,L)=P_{RR,X}(q_X(c_1(L))).
\end{equation}
This polynomial has the following properties:
\begin{enumerate}[{\rm (a)}]
\item\label{ita} the  constant   term of $P_{RR,X}(T)$ is $\chi(X,\cO_X)=n+1$;
\source{riemann-roch-polynomials-hyperkahler-fourfolds}
\item\label{itb} the leading term of $P_{RR,X}(T)$ is $\frac{c_X}{(2n)!} T^n $;
\source{riemann-roch-polynomials-hyperkahler-fourfolds}
\item the coefficients of $P_{RR,X}(T)$ are all positive (\cite[Theorem~1.1]{jia}).
\end{enumerate}

The next two lemmas are elementary.

\begin{lemm}\label{lem11}
\source{riemann-roch-polynomials-hyperkahler-fourfolds}
Let $X$ be a \hk\ manifold.\ For every class $\alpha\in H^2(X,\Z)$, one has $P_{RR,X}(q_X(\alpha))\in\Z$.
\end{lemm}

\begin{proof}
Since the period map is surjective (\cite[Proposition~25.12]{huy}), the class $\alpha$ becomes the first Chern class of some holomorphic line bundle on a deformation $X'$ of $X$, and~\eqref{rr} implies $P_{RR,X'}(q_{X'}(\alpha))\in \Z$.\ Since the polynomial~$P_{RR,X} $ and the form $q_X$ are deformation invariant, the lemma follows.
\end{proof}

\begin{lemm}\label{lemmaa}
\source{riemann-roch-polynomials-hyperkahler-fourfolds}
Let $X$ be a \hk\ manifold.\ Assume   that there is a  class~$\lll \in H^2(X,\Z)$ such that $\int_X \lll^{2n}=0$.\ For every $ \mm\in H^2(X,\Z)$, the number
\begin{equation}\label{defa2}
\source{riemann-roch-polynomials-hyperkahler-fourfolds}
a\coloneqq \frac{1}{n!}\int_X \lll^{n}\mm^{n}
\end{equation}
is  an integer.
\end{lemm}

\begin{proof}
By~\eqref{fuj}, we have $q_X(\lll)=0$.\ Set
\begin{equation}\label{defP}
\source{riemann-roch-polynomials-hyperkahler-fourfolds}
\forall k\in \Z\qquad P(k)\coloneqq P_{RR,X}(q_X(k\lll+ \mm))=P_{RR,X}(2kq_X(\lll,\mm)+q_X(\mm)) .
\end{equation}
Then $P$ is a  polynomial of degree $n$ whose leading coefficient is  (use~\eqref{pol} and property~\ref{itb} above)
\begin{equation}\label{defia}
\source{riemann-roch-polynomials-hyperkahler-fourfolds}
\frac{c_X}{(2n)!} (2 q_X(\lll,\mm))^n=\frac{1}{(n!)^2}\int_X \lll^{n}\mm^{n}=\frac{a}{n!}.
\end{equation}
By Lemma~\ref{lem11}, the polynomial $P$ takes integral values on integers, hence $a$ is  an integer.
\end{proof}


\begin{rema}
\source{riemann-roch-polynomials-hyperkahler-fourfolds}
\notready\label{prim}
\source{riemann-roch-polynomials-hyperkahler-fourfolds}
Under the hypotheses of Lemma~\ref{lemmaa}, when $a$ is    divisible by no nontrivial $n$th powers, the sublattice $\Z\lll\oplus \Z\mm$ of $H^2(X,\Z)$ is saturated: this follows from the fact that  $\frac{1}{n!}\int_X \lll^{n}\alpha^{n}$ is an integer for all $\alpha\in H^2(X,\Z)$.
\end{rema}




%%%%%%%%%%%%%%%%%%%%%
